\documentclass[10pt,a4paper]{amsart}
\usepackage[margin=19mm]{geometry}
\usepackage{amsmath,amssymb,mathtools}
\usepackage[scr=rsfs,cal=euler]{mathalfa}
\usepackage{xfrac}
\usepackage[unicode,hidelinks,bookmarks=false]{hyperref}
\newcommand{\ie}{i.e.\ }
\newcommand{\cf}{cf.\ }
\newcommand{\resp}{respectively\ }
\newcommand{\Subsection}{\subsection}
\providecommand{\nobreakdash}{\nobreak\hbox{-}}
\setlength{\emergencystretch}{2em}
\multlinegap0pt
\usepackage{bm}
\usepackage{bbm}
\usepackage{dsfont}
\usepackage{xspace}
\usepackage{cancel}
\usepackage{mathtools}
\usepackage{amsthm}
\makeatletter
\@ifundefined{theorem}{\newtheorem{theorem}{Theorem}}{}
\@ifundefined{lemma}{\newtheorem{lemma}[theorem]{Lemma}}{}
\makeatother
\DeclareMathOperator{\cp}{\,\square\,}
\DeclareMathOperator{\gp}{gp}
\title[Moving through Cartesian products, coronas and joins in gene]{Moving through Cartesian products, coronas and joins in general position}
\author{Sandi Klav\v{z}ar, Aditi Krishnakumar, Dorota Kuziak, Ethan Shallcross, James Tuite, Ismael G. Yero}
\date{Source: arXiv:2505.00535v2 (CC BY 4.0)}
\begin{document}
\maketitle

\noindent\emph{Source and licence.} Moving through Cartesian products, coronas and joins in general position. Version arXiv:2505.00535v2 (CC BY 4.0), distributed under the Creative Commons Attribution 4.0 International licence, as recorded in the arXiv licence metadata for this exact version and shown on the abstract page https://arxiv.org/abs/2505.00535v2. Full rights notice: licenses/arxiv-2505.00535-CC-BY-4.0.txt.\par\smallskip

\noindent\textit{Editorial scope.} Counted unit: the Lemma whose source label is \texttt{lem:unique-gp-in-Hamming} in the article (source0.tex lines 348--352, source label \texttt{lem:unique-gp-in-Hamming}), delivered together with the article's own complete original proof of it (source0.tex lines 354--371, one \texttt{proof} environment of 18 lines). No mathematics is written, repaired or re-derived: statement and proof are the article's own text line for line. Required context retained uncounted: none. Every definition, display and label of those lines is retained unchanged. Omitted from this package and not redistributed: 1--347, 353--353, 372--1447. Nothing inside a retained excerpt is omitted or altered: every retained body is delivered exactly as the article's lines are written, so no omission receipt of any kind accompanies this package. Citations: the retained excerpts carry no citation command at all, so no bibliography entry of the article is transcribed and no reference is redistributed. Reference adaptation: none was needed and none was made; every reference inside the delivered unit resolves inside it, so no unresolved reference or citation is produced. Printed slips kept literal: no printed word, symbol, hypothesis or number of the counted statement or of its proof was altered. Layout and class: the article's own class is \texttt{article} with packages including amssymb, amsthm, tikz, pgfplots, amsmath, todonotes, url, color, graphicx, placeins, caption, hyperref, geometry; the wrapper is the lane's pinned \texttt{amsart} editorial wrapper at 10pt on A4, which restates the theorem-like environments the retained text needs and adds the article's own macro definitions that the retained text needs (\texttt{\textbackslash cp}, \texttt{\textbackslash gp}), with extra wrapper packages (bm, bbm, dsfont, xspace, cancel, mathtools). The delivered numbering is the unit's own. Assets: no retained span contains a figure environment, a graphics-inclusion command, a PDF-page inclusion, a table or a TikZ picture; no image file is copied into this package, the package declares no asset dependency and no image file is redistributed with it. Licence: version arXiv:2505.00535v2 of this article is distributed under the Creative Commons Attribution 4.0 International licence, as recorded in the arXiv licence metadata for this exact version and shown on the abstract page https://arxiv.org/abs/2505.00535v2. A full rights notice is delivered at licenses/arxiv-2505.00535-CC-BY-4.0.txt. Characters: every retained excerpt is pure ASCII, so the delivered engine, XeLaTeX with its default Latin Modern text font, reports no missing character.
	\begin{lemma}
		\label{lem:unique-gp-in-Hamming}
		Let $n\ge 3$, $m\ge 3$, $V(K_n) = [n]$, $V(K_m) = [m]$, and let $X$ be a gp-set of $K_n\cp K_m$. Then there exist $i\in [n]$ and $j\in [m]$ such that 
		$$X = \left( V((K_n)^j) \cup V(^i(K_m))\right )\setminus \{(i,j)\}\,.$$ 
	\end{lemma}

	\begin{proof}
		Let $X$ be a gp-set of $K_n\cp K_m$. Then, as stated above, $|X| = n + m - 2$. Note that the set $X$ contains at most two vertices from every induced copy of $C_4$ of $K_n\cp K_m$. This fact implies the following: 
		
		\medskip\noindent
		{\bf Claim A}: if $(i,j), (i,j')\in X$, where $j\ne j'$, then 
		$X \cap \left( V((K_n)^j) \cup V((K_n)^{j'})  \right) = \{(i,j), (i,j')\}$. 
		
		\medskip
		If $X$ contains all the vertices of some $K_n$-layer, then by Claim~A, $X$ contains no other vertices, implying that $|X| = n < n + m - 2$, which is not possible. Let $j\in [m]$ be such that $t = |X \cap V((K_n)^{j})|$ is as large as possible. Note that $t\ge 2$, for otherwise we would have $|X|\le m$. Moreover, by the above, $t \le n-1$. If $t = n-1$ and $i\in [n]$ is the index for which $(i,j)\notin X$, then using Claim~A again, we have that $X \subseteq V((K_n)^j) \cup V(^i(K_m))$. Since $|X| = n+m-2$, we conclude that $X$ has the required structure in this case. Suppose finally that $t=n-k\ge 2$, where $k\ge 2$. We may assume without loss of generality that $X \cap V((K_n)^j ) = \{(1,j), \ldots, (n-k,j)\}$. Then, using Claim~A once more, 
		$$X \cap \bigcup_{\ell = 1}^{n-k} V((K_n)^{\ell}) = \{(1,j), \ldots, (n-k,j)\}\,.$$
		Notice that the subgraph $H$ of $K_n\cp K_m$ induced by the vertex set 
		$$\{n-k+1, \ldots, n\} \times \left( \{1, \ldots, j-1\} \cup \{j+1, \ldots, m \} \right)$$
		is isomorphic to $K_{k} \cp K_{m-1}$. Since it is a convex subgraph of $K_n\cp K_m$, the intersection $X\cap V(H)$ is a general position set of $H$. Therefore, 
		$$|X \cap V(H)| \le \gp(K_{k} \cp K_{m-1}) = k + (m-1) - 2 = k + m - 3\,.$$ 
		By our assumption on $t$ we have $X\cap \{(n-k+1,j), \ldots, (n,j)\} = \emptyset$. Thus we can conclude that $|X| \le (n - k) + (k + m - 3) = n + m -3$, a contradiction. 
		
		We have thus proved that $|X| = n + m - 2$ holds only in the case when $t = n-1$ and $X$ has the structure as claimed. 
	\end{proof}

\end{document}
